% !TEX root = main.tex
\section*{Nombres adimensionnels}

Les limites du modèle présentées ci-dessus peuvent être analysées 
plus quantitativement à l'aide de nombres adimensionnels. Le nombre 
de Péclet, le nombre de Damköhler et la longueur de Debye adimensionnée 
permettent d'évaluer l'importance relative des principaux mécanismes de 
transport et de réaction, et de préciser le domaine de validité de plusieurs 
hypothèses simplificatrices.

\subsection*{Péclet}

Le nombre de Péclet compare la convection et la diffusion :

\begin{equation}
Pe = \frac{vl_e}{D_{\text{ref}}}
\end{equation}

En partant des équations~\eqref{eq:nernst-planck},~\eqref{eq:conservation}, et en posant :

\begin{equation*}
c^* = \frac{C_i}{C_{\text{ref}}}, \quad 
\varphi_l^* = \frac{\varphi_l F}{RT}, \quad 
\nabla^* = l_e\nabla, \quad
t^* = \frac{t D_{\text{ref}}}{l_e^2}
\end{equation*}

En introduisant le flux adimensionné

\[
N_i^*=\frac{N_i l_e}{D_iC_{\rm ref}},
\]

on obtient :

\begin{equation*}
N_i^* = -\nabla^* c^* - z_i c^* \nabla^* \varphi_l^* + c^* Pe
\end{equation*}

et l'équation de conservation adimensionnée :

\begin{equation*}
\frac{\partial c^*}{\partial t^*} + \nabla^* \cdot \left(-\nabla^* c^* - z_i c^* \nabla^* \varphi_l^* + c^* Pe\right) = 0
\end{equation*}
%Je devrais sûrement enlevé la migration car là elle apparait aussi ou changé une phrase sans doute

Le nombre de Péclet compare uniquement les contributions convective et diffusive, 
il ne permet pas de quantifier l'importance du terme de migration.

% ============================================================
\subsection*{Damköhler}

Le nombre de Damköhler compare la vitesse de réaction surfacique à la diffusion :

\begin{equation}
Da = \frac{j l_e}{F D_i C_{\text{ref}}}
\end{equation}

En adimensionnant la condition limite~\eqref{eq:faraday}, on obtient :
%Faudrait faire attention car Damkoler c'est sur toute la surface pas que sur avcat

\begin{equation*}
-N_i^* \cdot \mathbf{n} = Da \frac{\nu_i}{n_e}
\end{equation*}

le flux normal étant nul sur le reste de l'interface. Ainsi défini, $Da$ caractérise 
la compétition entre réaction interfaciale et diffusion à l'échelle de la surface 
catalytique. Il est étroitement lié au module de Thiele~\cite{wanPotentialDependentThiele} 
introduit précédemment, qui décrit la même compétition à l'échelle de l'électrode, 
la couverture catalytique $\theta$ intervenant alors via l'aire spécifique active 
$a_v^{\text{cat}} = \theta\, a_v$.

% ============================================================
\subsection*{Longueur de Debye}

La longueur de Debye compare l'épaisseur caractéristique de la double
couche électrique à la taille caractéristique d'un pore $\ell_p$~\cite{dickinsonElectroneutralityApproximationElectrochemistry2011} :

\begin{equation}
\bar{\delta} = \frac{\lambda_0}{\ell_p}
= \frac{1}{\ell_p}
\sqrt{\frac{\varepsilon RT}{2F^2 C_{\text{ref}}}}
\end{equation}

En partant de l'équation de Poisson~\eqref{eq:poisson}, on adimensionne :

\begin{equation}
\quad \bar{\delta}^2 \nab^{*2} \varphi_l^* = -\frac{1}{2}\sum_i z_i c_i^*
\end{equation}

Lorsque la longueur de Debye est très inférieure à la taille caractéristique des pores, 
les effets de la double couche électrique restent confinés au voisinage des interfaces. 
L'approximation d'électroneutralité, introduite dans l'équation~\eqref{eq:electroneutre} lors de la 
présentation du modèle, est alors justifiée.

\bigskip

Cette analyse adimensionnelle complète les développements précédents en précisant le domaine 
de validité des principales hypothèses du modèle. Elle permet d’identifier, selon les conditions
 opératoires, les mécanismes susceptibles de gouverner le fonctionnement de l’électrode poreuse. 
 Les nombres de Péclet et de Damköhler, ainsi que la longueur de Debye adimensionnée, constitueront ainsi
  une grille de lecture pour les simulations numériques à venir et pour l’évaluation du domaine de validité des modèles réduits.